\section{InfluxDB Version Selection}
InfluxDB v2 OSS is selected as the historian for this project. The justification centres on four properties that are directly required by the system design.

Flux query language and native task scheduling. InfluxDB v2 is the only version for which Flux is the primary, fully--supported query interface, and all dashboard and analysis queries are written in Flux. The original downsampling design relied on Flux scheduled tasks running within the database; as built, downsampling is performed by the ingestion service instead (Section~\ref{sec:downsampling}), but the integrated task scheduler remains available without an external Kapacitor instance, keeping the Docker Compose stack minimal and operationally simple, appropriate for a single--historian research testbed.

Telegraf influxdb\_v2 output plugin and HTTP write API. The JSON ingestion path is entirely configuration--driven via Telegraf's mqtt\_consumer input and influxdb\_v2 output plugin. This plugin (including the \texttt{json\_time\_key} and \texttt{json\_time\_format = "unix\_ns"} features central to this design) is a first--class, mature integration with InfluxDB v2, with extensive documentation and community usage. The full-rate path writes batched, gzip-compressed line protocol directly to the same v2 HTTP API (Section~\ref{sec:fullrate}).

Token-based access control. InfluxDB v2's token--based authentication model supports bucket--scoped write tokens, which is the basis of the planned security posture (Section~\ref{sec:security}): write-only tokens scoped to their buckets for Telegraf and \texttt{daq2-ingest}, and a separate read-only token for Grafana. This fine-grained token model is native to v2, although the scoped tokens have not yet been issued.

Grafana Flux data source. Grafana's native Flux data source for InfluxDB v2 is stable and well--supported, enabling the dashboard panels (Section~\ref{sec:grafana}) to query the historian directly in Flux.

Relationship to InfluxDB 3 Core. InfluxDB 3 Core reached general availability in April 2025.  However, InfluxDB 3 Core is not positioned as a direct replacement for InfluxDB OSS v2, and v2 is not planned to reach end--of--life. The v3 architecture is a significant departure — it replaces Flux with native SQL and InfluxQL, and integrates a Python processing engine in place of the Kapacitor/task model used in v2.  Migrating to v3 would require rewriting all Flux queries documented here and the provisioned dashboard, with no clear benefit at this testbed scale. InfluxDB v2 OSS therefore remains the appropriate and supported choice for a new deployment of this scope.

Docker image pinning note: The latest tag for InfluxDB Docker images will point to InfluxDB 3 Core from late May 2026 onward.  The Compose file (Section~\ref{sec:compose}) uses the explicit influxdb:2 tag to prevent an unintended upgrade during the thesis validation period. This must not be changed to latest.
